% Generated draft: requires mathematical and source review.
\documentclass[11pt]{article}
\usepackage[margin=25mm]{geometry}
\usepackage{fontspec}
\setmainfont{Latin Modern Roman}
\usepackage{amsmath,amssymb,mathtools,unicode-math}
\setmathfont{Latin Modern Math}
\usepackage{enumitem,needspace,xurl,hyperref}
\setlength{\parindent}{0pt}
\setlength{\parskip}{5pt}
\setlength{\emergencystretch}{3em}
\newcommand{\N}{\mathbb{N}}
\newcommand{\Z}{\mathbb{Z}}
\newcommand{\Q}{\mathbb{Q}}
\newcommand{\R}{\mathbb{R}}
\newcommand{\C}{\mathbb{C}}
\newcommand{\F}{\mathbb{F}}
\newcommand{\A}{\mathbb{A}}
\newcommand{\PP}{\mathbb{P}}
\newcommand{\E}{\mathbb{E}}
\newcommand{\eps}{\varepsilon}
\newcommand{\dd}{\,\mathrm{d}}
\begin{document}
\section*{UnsolvedMath ID: 30004185}
\textbf{Source problem code:} catalogue entry
\subsection{Definitions and mathematical statement}
A traveling heteroclinic wave is a solution represented by a profile translating at a constant speed and approaching distinct equilibrium states at its two ends. In scalar profile notation, the heteroclinic boundary conditions are
\[
\lim_{\xi\to-\infty}U(\xi)=a_-,\qquad
\lim_{\xi\to+\infty}U(\xi)=a_+,\qquad a_-\ne a_+,
\]
where \(U\colon\R\to\R\) is the profile and \(a_-,a_+\) are equilibrium states associated with the two wells. This notation explains the terminology; the supplied record does not specify the governing equation or whether this scalar profile description is the intended higher-dimensional formulation.

\textbf{Extension problem.} Can the one-dimensional Frenkel--Kontorova existence theorem for traveling heteroclinic waves in suitable smooth anharmonic two-well potentials be extended to higher spatial dimensions \(d\ge 2\)?

The word ``suitable'' retains an unresolved qualification from the source. The supplied record does not give the precise one-dimensional theorem, its potential hypotheses, its admissible wave speeds, or the higher-dimensional equations and wave geometry. Consequently, a fully quantified existence statement cannot be extracted from this record alone. In particular, no assumptions about equal well depths, nondegenerate minima, propagation direction, or nonzero speed are supplied here.
\subsection{Short English statement}
The problem asks for a higher-dimensional analogue of a one-dimensional Frenkel--Kontorova theorem producing traveling waves that connect two distinct equilibrium states in a suitable smooth anharmonic two-well potential. The equations, potential assumptions, speed restrictions, and intended higher-dimensional wave geometry remain unspecified.
\subsection{Sources}
\begin{itemize}
\item Local UnsolvedMath registry: supplied complete record for permanent ID \(30004185\), titled ``Higher-Dimensional Frenkel--Kontorova Traveling Waves,'' in the Oberwolfach Reports Open Problems collection. Exact source problem code: \texttt{catalogue entry}. The record gives proposed year 2019 and names no proposer.
\item Supplied source citation: \emph{Mathematical Theory of Water Waves} (2020), Oberwolfach Reports, pp.~1919--1979, DOI \texttt{10.4171/owr/2019/32}. Original report: \url{https://doi.org/10.4171/owr/2019/32}. The supplied background describes Oberwolfach Reports as documenting mathematical workshops and their open-problem discussions.
\item Supplied literature reference: Buffoni--Schwetlick--Zimmer, ``heteroclinic waves'' (the supplied reference label), \url{https://doi.org/10.1016/j.matpur.2019.01.002}.
\item Supplied literature assessment dated exactly \texttt{2026-08-22}: the higher-dimensional extension is posed as a question, and the cited primary work establishes only the one-dimensional-chain result. The record labels the problem open. This is supplied dated curation derived from the OWR material, not an independently verified literature assessment; the status and cited claims require independent verification.
\end{itemize}
\end{document}
